% Annexes du guide d'annotation
\section{Annexes}

\subsection{Glossaire}

\textbf{Bactériémie} : Présence de bactéries dans le sang, constituant une infection bactérienne systémique potentiellement grave.

\textbf{CIM-10} : Classification Internationale des Maladies, 10ème révision. Système de codification médicale utilisé pour classifier les maladies et les problèmes de santé.

\textbf{Score F1} : Mesure de performance combinant la précision et le rappel selon la formule harmonique : F1 = 2 × (précision × rappel) / (précision + rappel).

\textbf{Précision} : Proportion de vrais positifs parmi l'ensemble des éléments identifiés comme positifs par le système.

\textbf{Rappel} : Proportion de vrais positifs identifiés par le système parmi l'ensemble des vrais positifs existants.

\textbf{Site primaire} : Localisation anatomique d'origine de l'infection bactérienne qui a conduit à la bactériémie.

\textbf{Site secondaire} : Localisation anatomique où l'infection s'est propagée secondairement à partir du site primaire.

\textbf{Traitement automatique du langage (TAL)} : Ensemble de techniques informatiques permettant l'analyse, la compréhension et le traitement automatique de textes en langage naturel.

\textbf{Algorithme à base de règles} : Système de traitement automatique s'appuyant sur des règles linguistiques et des terminologies prédéfinies pour extraire des informations spécifiques.

\textbf{Terminologie de négation} : Ensemble de termes et expressions utilisés pour identifier les formulations négatives dans les textes médicaux.

\textbf{BERT} : Bidirectional Encoder Representations from Transformers. Architecture de modèle de langage basée sur les transformeurs.

\textbf{Modèles de langage de grande taille} : Modèles d'intelligence artificielle entraînés sur de vastes corpus textuels pour comprendre et générer du langage naturel.

\textbf{Approches supervisées} : Méthodes d'apprentissage automatique utilisant des données étiquetées pour entraîner des modèles prédictifs.

\textbf{BACTHUB} : Projet de l'Entrepôt de Données de Santé de l'AP-HP portant sur l'analyse des bactériémies à partir des dossiers hospitaliers.

\textbf{Projet PARTAGES} : Projet de recherche collaborative portant sur le développement et l'évaluation de modèles de traitement automatique du langage dans le domaine médical.